\chapter*{References}
\addcontentsline{toc}{chapter}{References}

\begin{list}{}{%
  \setlength{\leftmargin}{1.6em}%
  \setlength{\itemindent}{-1.6em}%
  \setlength{\itemsep}{0.55em}%
  \setlength{\parsep}{0pt}%
  \setlength{\topsep}{0.6em}%
}

\item\phantomsection\label{cite:baird1995}%
Baird, L. (1995). Residual algorithms: Reinforcement learning with function approximation. In \textit{Proceedings of the 12th International Conference on Machine Learning} (pp.\ 30--37).

\item\phantomsection\label{cite:barto1983}%
Barto, A. G., Sutton, R. S., and Anderson, C. W. (1983). Neuronlike adaptive elements that can solve difficult learning control problems. \textit{IEEE Transactions on Systems, Man, and Cybernetics}, \textit{SMC-13}(5), 834--846.

\item\phantomsection\label{cite:bellman1957}%
Bellman, R. (1957). \textit{Dynamic Programming}. Princeton University Press.

\item\phantomsection\label{cite:christiano2017}%
Christiano, P. F., Leike, J., Brown, T. B., Martic, M., Legg, S., and Amodei, D. (2017). Deep reinforcement learning from human preferences. In \textit{Advances in Neural Information Processing Systems 30} (pp.\ 4299--4307).

\item\phantomsection\label{cite:deepseek2025}%
Guo, D., Yang, D., Zhang, H., et al. (2025). DeepSeek-R1 incentivizes reasoning in LLMs through reinforcement learning. \textit{Nature}, \textit{645}(8081), 633--638.

\item\phantomsection\label{cite:jensen1906}%
Jensen, J. L. W. V. (1906). Sur les fonctions convexes et les inégalités entre les valeurs moyennes. \textit{Acta Mathematica}, \textit{30}, 175--193.

\item\phantomsection\label{cite:kool2019}%
Kool, W., van Hoof, H., and Welling, M. (2019). Buy 4 REINFORCE samples, get a baseline for free! In \textit{ICLR 2019 Workshop on Deep Reinforcement Learning Meets Structured Prediction}.

\item\phantomsection\label{cite:mnih2015}%
Mnih, V., Kavukcuoglu, K., Silver, D., Rusu, A. A., Veness, J., Bellemare, M. G., Graves, A., Riedmiller, M., Fidjeland, A. K., Ostrovski, G., Petersen, S., Beattie, C., Sadik, A., Antonoglou, I., King, H., Kumaran, D., Wierstra, D., Legg, S., and Hassabis, D. (2015). Human-level control through deep reinforcement learning. \textit{Nature}, \textit{518}(7540), 529--533.

\item\phantomsection\label{cite:mnih2016}%
Mnih, A. and Rezende, D. J. (2016). Variational inference for Monte Carlo objectives. In \textit{Proceedings of the 33rd International Conference on Machine Learning} (pp.\ 2188--2196).

\item\phantomsection\label{cite:ouyang2022}%
Ouyang, L., Wu, J., Jiang, X., Almeida, D., Wainwright, C. L., Mishkin, P., Zhang, C., Agarwal, S., Slama, K., Ray, A., Schulman, J., Hilton, J., Kelton, F., Miller, L., Simens, M., Askell, A., Welinder, P., Christiano, P., Leike, J., and Lowe, R. (2022). Training language models to follow instructions with human feedback. In \textit{Advances in Neural Information Processing Systems 35} (pp.\ 27730--27744).

\item\phantomsection\label{cite:robbins1951}%
Robbins, H. and Monro, S. (1951). A stochastic approximation method. \textit{The Annals of Mathematical Statistics}, \textit{22}(3), 400--407.

\item\phantomsection\label{cite:rummery1994}%
Rummery, G. A. and Niranjan, M. (1994). On-line Q-learning using connectionist systems (Technical Report CUED/F-INFENG/TR 166). Cambridge University Engineering Department.

\item\phantomsection\label{cite:schulman2017}%
Schulman, J., Wolski, F., Dhariwal, P., Radford, A., and Klimov, O. (2017). Proximal policy optimization algorithms. arXiv:1707.06347.

\item\phantomsection\label{cite:shao2024}%
Shao, Z., Wang, P., Zhu, Q., Xu, R., Song, J., Bi, X., Zhang, H., Zhang, M., Li, Y. K., Wu, Y., and Guo, D. (2024). DeepSeekMath: Pushing the limits of mathematical reasoning in open language models. arXiv:2402.03300.

\item\phantomsection\label{cite:stiennon2020}%
Stiennon, N., Ouyang, L., Wu, J., Ziegler, D. M., Lowe, R., Voss, C., Radford, A., Amodei, D., and Christiano, P. (2020). Learning to summarize from human feedback. In \textit{Advances in Neural Information Processing Systems 33}.

\item\phantomsection\label{cite:sutton1988}%
Sutton, R. S. (1988). Learning to predict by the methods of temporal differences. \textit{Machine Learning}, \textit{3}, 9--44.

\item\phantomsection\label{cite:suttonbarto2018}%
Sutton, R. S. and Barto, A. G. (2018). \textit{Reinforcement Learning: An Introduction} (2nd ed.). MIT Press.

\item\phantomsection\label{cite:sutton2000}%
Sutton, R. S., McAllester, D., Singh, S., and Mansour, Y. (2000). Policy gradient methods for reinforcement learning with function approximation. In \textit{Advances in Neural Information Processing Systems 12} (pp.\ 1057--1063).

\item\phantomsection\label{cite:tsitsiklis1997}%
Tsitsiklis, J. N. and Van Roy, B. (1997). An analysis of temporal-difference learning with function approximation. \textit{IEEE Transactions on Automatic Control}, \textit{42}(5), 674--690.

\item\phantomsection\label{cite:vanhasselt2010}%
van Hasselt, H. (2010). Double Q-learning. In \textit{Advances in Neural Information Processing Systems 23} (pp.\ 2613--2621).

\item\phantomsection\label{cite:vanhasselt2016}%
van Hasselt, H., Guez, A., and Silver, D. (2016). Deep reinforcement learning with double Q-learning. In \textit{Proceedings of the AAAI Conference on Artificial Intelligence} (pp.\ 2094--2100).

\item\phantomsection\label{cite:watkins1992}%
Watkins, C. J. C. H. and Dayan, P. (1992). Q-learning. \textit{Machine Learning}, \textit{8}, 279--292.

\item\phantomsection\label{cite:williams1992}%
Williams, R. J. (1992). Simple statistical gradient-following algorithms for connectionist reinforcement learning. \textit{Machine Learning}, \textit{8}, 229--256.

\item\phantomsection\label{cite:ziegler2019}%
Ziegler, D. M., Stiennon, N., Wu, J., Brown, T. B., Radford, A., Amodei, D., Christiano, P., and Irving, G. (2019). Fine-tuning language models from human preferences. arXiv:1909.08593.

\end{list}
